\documentclass{jsarticle}
\usepackage{amsmath}
\usepackage{amssymb}
\begin{document}
\title{}
\author{}
\date{}
\maketitle

  Pholographic reading systemを、今までの速読法から作り上げる。
まずは、SRS速読法を気功法からこのPholographic reading systemに取り上げる。
凝視は、ページをチュッチュチュッチュ文字を食らいつくのに使う。
この気功法は、瞑想法も入っているので、潜在意識に情報を流し入れられる。
統合出力のときも、この瞑想法で、潜在意識に入れている情報を汲み取るのも
自然に出来る。次に、Photo reading hole mind systemも採り入れる。
みかん集中法は、視界が開けて、その上に、SRS速読法の凝視の疲れもとれる。
逆さ読みは、メガネを付けなかったら、凝視が入るが、メガネを付けるとこの疲れも
精神には入らない。アフォメーションは、始めと終わりに儀式としていれると、
目的に可能性から無理がないと、思いが実現する。
SRSの指回し体操は、潜在意識に入っている夢を想起するのに使う。SRS気功法
も、潜在意識に入っている夢の想起にも使える。
Photo reading hole mind systemの快楽も夢実現にも使える。SRS速読法は、
指回し体操がドーパミン遮断なのに快楽が出て来る。この指回し体操は、
速めくりと共鳴力と気功法がすべてセットになっている。
GSR速読法のGenerative state zenは、頭頂葉から下半身に気の巡りを通らして、
精神のリラックスをつくり、楽に潜在意識に情報を入力する役割がある。
リモートアイは、モンロー研究所が調べている。イメージトレーニング
に、１２種類の眼球の筋肉を鍛える機構が入っている。
瞬読は、ゲシュタルトトレーニングに単語の変換力を使って、速読を
バラバラの情報から読み取るのに使える。
散歩法は、この家で訓練しているときの潜在意識に入れ込んだ情報を、
消化して吸収する訓練になっている。
蔵王病院で、最近になって、初めて診察を待つときに瞑想していたときに、
頭の栓が抜けた感じで、人の話し声が、素通りして、その上に快楽で、
精神がリラックスしていた。この瞑想法はPholographic reading systemに使える。
栗田博士は、一日の疲れの時に、バルコニーに座って、一日の思いをだんだんと
意識のレベルを落して、この快楽で、潜在意識の中に入り込むらしい。
日記に書く時に、前頭葉を全開にして、潜在意識に入り込んでいる情報を
音読でトリガーとアンカーの技法を使って書きまくる。
私が、初めて論文を書くのに、彩さんの一冊目の本の自分で分かった言葉と
読んだ内容を自分の言葉で書くと、統合出力できますよが、速読法の
統合出力のきっかけになった。今までは、読んだ内容そのままで書こうとして
いたから、統合出力が出来ていなかった。
河相先生のレボトミンは、リボ核酸だけでなく、側頭葉の精神のバランスが
くずれたためのドーパミンの量を調節するのに、やさしく側頭葉をブロックする。
SRS速読法もこの機構を使っている。そのために、側頭葉を活性化させずに、
情報の処理を頭頂葉にもっていっている。統合出力のときに、側頭葉を全開に使う。
内容の統合をするから、書く時は、側頭葉を活性化させてもいいらしい。
萬谷先生は、超能力の薬を出していた。セレンとテルルは、ウィスパードの
サナギで調べられている。
アカシックレコードは、前世療法と未来世療法で見に付く。
音読の良さは、声が綺麗になることと、統合出力の時に役立つ事だ。
それゆえに、速読法と音読法は、18年間隔で流行が来る。
勝間さんは、ビデオテープのように、一日の情報を早送りして、見ている。

リスペリドンは、側頭葉の緩和と精神の安定によい。エヴィリファイと違って、
過鎮静がない。周りが恐くなる事もない。すごく屯服を服用していたら、ものすごく
楽でよかった。これ以上の薬は、ロナセンの規定量しかないとおもう。
睡眠薬で、平日に飲んでいいエチゾラムは、瞑想状態に近かった。

速読法は、気は血に乗って動くというのと、共鳴力で周りと情報を共有するらしい。
牽引の法則で、先を見通すこともあるらしい。アカシックレコードは、応用すると
これらの現象が現れるらしい。

機知を司る第32野と第33野は、予知も入っているらしい。それゆえに、安岡先生は
五行を捨てないようにと言っている。劫敗占星型と劫敗食傷型が、私の能力らしい。
真逆は、それでも占いは出来づらいらしい。みんなは見えているらしい。
酉の能力は、契約と実現らしい。星座占いは、9月と6月が反対らしい。

今までの政治や宗教、戦争の時期は、各国の指導者が易学で周りに悟られずに
実行していたらしい。史学を年号で卦を出すと最近分かった。

音読は、一番周りにも自分にもやさしい。

これらの速読法と音読法を統合して、Pholographic reading systemをつくっていく。

まとめると、イメージトレーニングの瞑想法と退行催眠と未来世療法のアカシック
レコード、気功法と統合出力、変換イメージトレーニングのゲシュタルト訓練を
統合して、Pholographic reading systemを形成していく。 

これらの速読法と音読法を合わせると、勝間さんのビデオの早送りと巻き戻しに
なるらしい。



\end{document}
